\subsection{Finite saddle splittings and the two center axes}

The two lower contacts at a finite upper saddle require independent
splitting coordinates, including near the intersection of the two
center families. Fix $\bar B=-\beta$ with
$0<\beta_*\le\beta\le\beta^*$, and work near
$\bar m=\beta$, $\bar c=-\bar a$, $\nu=0$.
The lower roots are simple on this compact block. Let
$\widehat p_+$ and $-\widehat p_-$ be their strong intersections with
$Y=0$, and let $P_s$ and $-Q_s$ be the intersections of the upper
saddle's two separatrices with that same section. Set
\[
 \delta_+=\widehat p_+-P_s,\qquad
 \delta_-=\widehat p_--Q_s.
\]
Each intersection is defined by a single transverse connecting arc.
We can therefore define both splittings before forming the complete
upper passage; the following Jacobian calculation makes them
coefficient coordinates.

\begin{lemma}[Finite splitting coordinates]
\label{rw:coordinates:finite-splittings}
The variables $(\delta_+,\delta_-,r,\nu)$, with $\beta$ fixed,
are coefficient coordinates on a common neighborhood of each compact
part of the three-saddle center, where $r$ is the actual upper
eigenvalue ratio. Near the reversible intersection, choose reflected
lower power inputs and put
\begin{equation}\label{rw:coordinates:finite-axes}
 u_0=\chi_+(\delta_+),\quad z_0=\chi_-(\delta_-),\qquad
 \xi_1=(u_0+z_0)/2,\quad \xi_2=(u_0-z_0)/2,\quad \kappa=r-1.
\end{equation}
Then $(\xi_1,\xi_2,\kappa,\nu)$ are coordinates with bounded
finite inverse jets on a common product box. The two center components
in this box are exactly
\begin{equation}\label{rw:coordinates:finite-center-axes}
 \xi_1=\xi_2=\nu=0,
 \qquad\text{and}\qquad \xi_2=\kappa=\nu=0.
\end{equation}
At fixed $\kappa$ near their intersection,
$\partial_{\bar a}\xi_2>c>0$ after the other independent coefficients
are held fixed. These statements use ordinary coefficient and Euler
scale derivatives of any finite order reserved in the preparations.
\end{lemma}

\begin{proof}
The lower strong intersections and positive maps $\chi_\pm$ are
supplied by Proposition~\ref{rw:lower-preparation}. The upper axes
and their outer connections are supplied by
Section~\ref{rw:upper-resonance}. For the present coordinate assertion
only its invariant-axis construction is needed: the denominators
$n\lambda_u-\lambda_s$ and $\lambda_u-n\lambda_s$ have the same
positive linear margins on every fixed compact set of positive saddle
ratios. Thus the axes and regular legs apply on the stated compact
center block without the later resonant shears. The compensator
estimates remain restricted to the neighborhood of ratio one.
Continuation to $Y=0$ uses the low-layer carrier estimate, which gives
convergence of the coefficient jets as $\epsilon\to0$. Since $|X|$
has a positive lower bound on the final low layer, the connection
remains uniformly transverse. We compute the coordinate Jacobian at
zero scale and preserve its margin by one common reduction of the
scale and coefficient domains.

Write $\Delta=\sqrt{\bar a^2+4\beta}$. The zero-scale lower
roots and upper thresholds are
\begin{equation}\label{rw:coordinates:finite-limit}
\begin{aligned}
 p_l&=\frac{\bar a+\Delta}{2\beta},&
 q_l&=\frac{-\bar a+\Delta}{2\beta},\\
 P_s&=\frac2{\bar c+\sqrt{\bar c^2+4\bar m}},&
 Q_s&=\frac2{-\bar c+\sqrt{\bar c^2+4\bar m}}.
\end{aligned}
\end{equation}
The upper formulas follow by intersecting the affine saddle's two
eigenlines with $Y=0$. On the center put $p=p_l=P_s$, $q=q_l=Q_s$.
Implicit differentiation of
$\bar mP_s^2+\bar cP_s=1$ and
$\bar mQ_s^2-\bar cQ_s=1$ gives
\[
 \frac{\partial(\delta_+,\delta_-)}{\partial(\bar m,\bar c)}
 =\frac1\Delta
   \begin{pmatrix}p^2&p\\q^2&-q\end{pmatrix},\qquad
 \det=-\frac1{\beta^2\Delta}.
\]
On the zero-splitting axis, $\bar m=\beta$, $\bar c=-\bar a$,
the ratio is $r=(\Delta+\bar a)/(\Delta-\bar a)$, with
$\dd r/\dd\bar a=2r/\Delta>0$. The two splitting directions
and this remaining tangent direction therefore give an invertible
three-variable map. The exact triangle identity of
Lemma~\ref{rw:centers:triangle} supplies the positive-scale
zero-splitting axis.

At $\bar a=\bar c=0$, choose the two reflected power inputs to
have the same positive linear coefficient $k_l(\beta)$ at zero.
The Jacobian in $(\bar m,\bar a,\bar c)$ is then
\begin{equation}\label{rw:coordinates:finite-jacobian}
 \frac{\partial(\xi_1,\xi_2,\kappa)}
      {\partial(\bar m,\bar a,\bar c)}
 =\begin{pmatrix}
 k_l/(2\beta^{3/2})&0&0\\
 0&k_l/(2\beta)&k_l/(2\beta)\\
 0&0&-1/\sqrt\beta
 \end{pmatrix},\qquad
 \det=-\frac{k_l^2}{4\beta^3}.
\end{equation}
Coefficient-jet convergence preserves these margins, so
Lemma~\ref{rw:coordinates:inverse-box}, with $\beta$ and scale
passive, supplies the common inverse boxes and jets.
The positive $\bar a$ derivative at fixed $\kappa$ follows from
the middle diagonal entry and continuity.

On the second center both separatrices are the triangle sides, so
$u_0=z_0=0$. Conversely the inverse coordinates and the nonzero tangent
ratio derivative identify this entire nearby axis with that component.
On the reversible component, reflection gives $u_0=z_0$ and $r=1$.
Along it $\partial_{\bar m}\xi_1$ has the positive limit in
\eqref{rw:coordinates:finite-jacobian}; it parametrizes the whole
second axis in \eqref{rw:coordinates:finite-center-axes} on a smaller
box. Thus the two center components fill precisely the two axes in
\eqref{rw:coordinates:finite-center-axes}.
\end{proof}

Division of the upper orbital remainder requires one coordinate choice
whose remainder vanishes on both center components. We now construct
that choice together with its regular outer legs.

\begin{lemma}[One upper coordinate choice on both centers]
\label{rw:coordinates:two-center-upper}
Near the intersection above, the analytic upper coordinates can be
chosen so that their orbital equation is
\[
 \frac{\dd V}{\dd U}=-r\frac VU-V^2R(U,V),\qquad
 R=\epsilon\widetilde R,
\]
and $R=0$ identically on each of the two center components. Their
finite jets and inverse jets have the same uniformity as the finite
saddle preparation. The two regular outer legs, expressed in the lower
power inputs, have the form
\[
 G_+(t)=u_0+tA_0(t),\qquad
 G_-(v)=z_0+vB_0(v),\qquad G_+',G_-'>c>0,
\]
on fixed signed connection intervals and on the common coefficient box.
\end{lemma}

\begin{proof}
On the triangle use the coordinates from
Lemma~\ref{rw:centers:triangle}; its first integral is
$U^{\lambda_+}V^{\lambda_-}$, so its orbital equation is exactly
$V_U=-rV/U$. On the reversible component the product coordinates
constructed in Section~\ref{rw:upper-resonance} give $V_U=-V/U$.
At their intersection $r=1$. The transition between these two analytic
coordinates preserves the foliation $UV=\text{constant}$. To see the
regularity at the axes, express either analytic first integral in the
other product coordinates. Invariance under
$U\partial_U-V\partial_V$ makes its power series a function of
$UV$ alone. Its derivative in $UV$ is nonzero, by the saddle Hessian.
Its analytic one-variable inverse shows that the transition and its
inverse preserve the product foliation on a whole saddle box,
including the axes.

Apply this intersection transition, with its passive $\beta,\epsilon$
dependence, to the entire reversible family. Its orbital equation stays
exactly linear and its coordinates now agree with the triangle choice
at the intersection. In coefficient variables
$\tau=\bar a+\bar c$, $\eta=\bar m-\beta$, extend the two
coordinate maps by taking their sum minus their common restriction:
the triangle prescription is evaluated at $\eta=\tau=\nu=0$,
the reversible prescription at $\bar a=\bar c=\nu=0$, and the
subtracted intersection prescription at all four zero values.
Use $\bar a$ as the remaining coordinate on the triangle and $\eta$
on the reversible family, then extend independently in $\tau,\nu$.
The maps agree on the intersection and retain an invertible state
Jacobian on a smaller common neighborhood. Subtract the image of the
actual saddle to center it; this subtraction vanishes on both prescribed
families.

Straighten the neighboring invariant axes and normalize the orbital
coefficient on them, choosing every change to be the identity when the
axes and coefficient are already normalized. The two center equations
are preserved. At zero scale both prescriptions and the field are
affine-linear, so the normalized remainder vanishes there; division
by $\epsilon$ gives its bounded analytic quotient. The shear
preparation and reciprocal estimates of
Proposition~\ref{rw:upper:prepared-map} apply with the same finite
ratio restriction, now preserving both centers.

Each outer leg to $Y=0$ lies in a compact regular strip except for
its final low layer, whose uniform jets were proved in Section~4.
Compose it with the regular connection into the lower power input.
The orientation is inward from the corresponding strong barrier, so
the resulting derivatives $G_+',G_-'$ are positive with fixed margins.
Reflect the sections and coordinate prescriptions together to obtain
the two legs. The signed intervals extend these individual regular
connections; the upper passage retains its original saddle quadrant.
Choose $|u_0|,|z_0|$ smaller than a fixed fraction of the signed
connection-image radius. The common inverse-box argument then gives
the inverse jets on a smaller image.
\end{proof}
